\section{What The Foundation Carries Forward}

The chapter has built only what the grammar forced it to build. It
began with distinguishability because sameness was too strong to be
primitive. It turned distinction into counting by showing that a reading
can preserve a whole or expose a part. It treated real value as finite
approach with residue, not as a completed infinity. It introduced time
as alternation and measurement as repeatable trial. It accepted a
halting horizon because a record must stop in order to be read.

Then the grammar pointed the tower at truth. At the bare foundation,
affirmation and denial could not be internally separated. The collapse
was not a denial of ordinary truth. It was the discovery that ordinary
truth depends on a standing place the foundation had not yet earned. The
needle supplied that place by one bounded act of manufactured identity.

What passes into the next chapter is therefore a disciplined grammar of
measurement. It permits difference before identity. It forbids completed
infinity from masquerading as a finite record. It carries residue as the
honest remainder of approximation. It bounds observation by trial time.
It quotients only where identity has been sanctioned. It restores truth
without making truth a primitive exemption from measurement.

The finite gauge will use exactly these inheritances. It will ask how a
path can be read when only finite conditions are available, how an
action can carry residue, and how stationarity can be identified without
abandoning the bound. The first chapter gives it the place to stand: a
needle after collapse, and a grammar that knows the cost of every
identity it uses.
